\section{Related Work}
\label{sec:related}

\paragraph{Structured dynamics and Hamiltonian learning.}
Hamiltonian and port-Hamiltonian learning provide the closest mechanistic
precedent for our view of policy as a structured phase-space flow.
Hamiltonian neural networks learn a scalar energy whose symplectic
gradient defines the vector field~\citep{greydanus2019hamiltonian};
SymODEN extends this idea to controlled systems~\citep{zhong2020symplectic};
Lagrangian networks and port-Hamiltonian neural networks add coordinate
or dissipation structure for physical system identification and
control~\citep{cranmer2020lagrangian,desai2021porthamiltonian}.
Classical port-Hamiltonian theory and energy-shaping control likewise
show how damping, interconnection, and external ports can be used to
shape closed-loop behavior~\citep{vanderschaft2014port}.
These works motivate our parameterization, but their usual target is
dynamics identification or stabilizing control from full-state data.
In contrast, we use the Hamiltonian structure as a navigation policy
class and add material-risk terms that act directly as learned force
channels under local sensing.

\paragraph{Safe and risk-sensitive reinforcement learning.}
Safe RL typically introduces risk as an objective- or constraint-level
quantity. CVaR optimization provides the variational foundation for
tail-risk objectives~\citep{rockafellar2000optimization}; CVaR MDPs and
risk-constrained actor-critic methods optimize tail or percentile risk
over trajectories~\citep{chow2015risk,chow2017risk}; and constrained
policy optimization enforces CMDP-style cost constraints during policy
search~\citep{achiam2017constrained}. Distributional RL instead models
the return distribution and can apply risk distortions to that
distribution~\citep{bellemare2017distributional,dabney2018implicit}.
Our objective uses the same tail-risk lineage, but the object being
updated is different: the CVaR signal trains coefficients of a structured
decision field, so risk changes the generated forces rather than only a
generic policy parameter vector, Lagrange multiplier, or return statistic.

\paragraph{Navigation, local planning, and safety filters.}
Classical navigation methods adapt paths or local actions through
search, sampling, or hand-designed potentials. A* and sampling-based
planners provide strong geometric references when map access is
available~\citep{hart1968formal,karaman2011sampling}; artificial
potential fields, the Dynamic Window Approach, MPPI, and CBF-QP filters
produce reactive local behavior under limited information
~\citep{khatib1986real,fox1997dynamic,williams2017mppi,ames2016control}.
These systems are important baselines because they also operate through
local costs or forces, but their risk terms are usually hand-designed,
externally scored, or enforced as filters. Our method keeps the local
closed-loop character while learning how semantic/material risk enters
the phase-space update.

\paragraph{Material-aware perception and traversability.}
Terrain and traversability methods often learn semantic classes,
freespace, or cost maps that are then consumed by a planner or
controller. Off-road datasets such as RUGD and RELLIS-3D support this
line of work~\citep{rugd2019,rellis2020}, and systems such as TerraPN
learn terrain costs for downstream local planning~\citep{sathyamoorthy2022terrapn}.
The DFC2018 Houston data provide a complementary remote-sensing setting
with rich land-cover cues~\citep{dfc2018}. Our DFC experiments use these
material labels to construct risk fields, but the contribution is not a
new semantic mapper: it is a mechanism for turning material risk into
internal energy and force terms that reshape the decision field.

\paragraph{Continual and multi-timescale adaptation.}
A final related thread studies adaptation over multiple timescales:
actor-critic methods separate critic and actor updates, continual
learning protects or reuses knowledge across tasks, and rapid motor
adaptation conditions policies on recent history
~\citep{kirkpatrick2017ewc,rusu2016progressive,kumar2021rma}.
Our three loops have a similar motivation--fast local correction,
episodic tail-risk updates, and slower curriculum advancement--but the
state being adapted is explicitly structured: local corrections and
episodic gradients modify Hamiltonian coefficients and material force
channels rather than an unconstrained black-box policy.

Prior work therefore covers each ingredient in isolation: structured
energy-based dynamics, risk-sensitive objectives, local planners and
safety filters, traversability maps, and multi-timescale adaptation. Our
contribution is their combination at the point where risk enters the
closed loop: a material-energy enrichment that reshapes the Hamiltonian
decision field itself.
